\usepackage[T1]{fontenc}
\usepackage[utf8]{inputenc}
\usepackage[spanish,es-nodecimaldot]{babel}
\usepackage{lmodern}

\usepackage{microtype}
\usepackage[a4paper,margin=2.8cm]{geometry}
\usepackage{setspace}
\onehalfspacing

\usepackage{amsmath,amssymb,amsfonts,mathtools}
\usepackage{bm}
\usepackage{siunitx}
\sisetup{detect-all}

\usepackage{graphicx}
\usepackage{booktabs}
\usepackage{tabularx}
\usepackage{longtable}
\usepackage{ragged2e}
\usepackage{caption}
\usepackage{float}
\usepackage{placeins}
\usepackage{adjustbox}
\usepackage{multirow}
\usepackage{array}
\newcolumntype{L}{>{\RaggedRight\arraybackslash}X}
\newcolumntype{C}{>{\Centering\arraybackslash}X}
\newcolumntype{P}[1]{>{\RaggedRight\arraybackslash}p{#1}}

\usepackage{enumitem}
\usepackage{csquotes}
\usepackage[table]{xcolor}
\usepackage[numbers,sort&compress]{natbib}
\usepackage[colorlinks=true,allcolors=blue]{hyperref}
\usepackage{xurl}
\urlstyle{tt}
\usepackage[spanish,nameinlink,noabbrev]{cleveref}
\crefname{chapter}{capítulo}{capítulos}
\Crefname{chapter}{Capítulo}{Capítulos}
\crefname{section}{sección}{secciones}
\Crefname{section}{Sección}{Secciones}
\crefname{subsection}{sección}{secciones}
\Crefname{subsection}{Sección}{Secciones}
\crefname{equation}{ecuación}{ecuaciones}
\Crefname{equation}{Ecuación}{Ecuaciones}
\crefname{table}{cuadro}{cuadros}
\Crefname{table}{Cuadro}{Cuadros}
\crefname{figure}{figura}{figuras}
\Crefname{figure}{Figura}{Figuras}
\usepackage{tikz}
\usepackage[titles]{tocloft}
% Espacio para números como 12.14.3 en el índice del manuscrito ampliado.
\setlength{\cftchapnumwidth}{2.3em}
\setlength{\cftsecindent}{2.3em}
\setlength{\cftsecnumwidth}{3.3em}
\setlength{\cftsubsecindent}{5.6em}
\setlength{\cftsubsecnumwidth}{4.2em}

\usetikzlibrary{arrows.meta,positioning,fit,shapes.geometric,babel}

\definecolor{azulacademico}{RGB}{45,73,104}
\definecolor{azulclaro}{RGB}{220,229,238}
\definecolor{griscronograma}{RGB}{225,225,225}

% ======================================================
%                   ACRÓNIMOS
% ======================================================
\usepackage[acronym,nogroupskip]{glossaries-extra}
\setabbreviationstyle[acronym]{long-short}
\makeglossaries
\newcommand{\ac}{\gls}
\newcommand{\acp}{\glspl}
\newcommand{\acs}{\glsentryshort}
\newcommand{\acf}{\glsfirst}

% RF / Comms / Sensing
\newacronym{RF}{RF}{radiofrecuencia}
\newacronym{WiFi}{Wi-Fi}{tecnología de redes inalámbricas basada en IEEE 802.11}
\newacronym{LoRa}{LoRa}{radio de largo alcance}
\newacronym{LoRaWAN}{LoRaWAN}{red de área amplia y largo alcance}
\newacronym{LPWAN}{LPWAN}{red de área amplia y bajo consumo}
\newacronym{OFDM}{OFDM}{multiplexación por división de frecuencias ortogonales}
\newacronym{CSI}{CSI}{información del estado del canal}
\newacronym{CFR}{CFR}{respuesta en frecuencia del canal}
\newacronym{CIR}{CIR}{respuesta al impulso del canal}
\newacronym{MIMO}{MIMO}{múltiples entradas y múltiples salidas}
\newacronym{SISO}{SISO}{una entrada y una salida}
\newacronym{RT}{RT}{trazado de rayos}
\newacronym{DRT}{DRT}{trazado de rayos diferenciable}
\newacronym{WDT}{WDT}{gemelo digital inalámbrico}
\newacronym{ISAC}{ISAC}{comunicación y sensado integrados}
\newacronym{SDR}{SDR}{radio definida por software}
\newacronym{IQ}{I/Q}{componentes en fase y cuadratura}
\newacronym{AoA}{AoA}{ángulo de llegada}
\newacronym{AoD}{AoD}{ángulo de salida}
\newacronym{ToF}{ToF}{tiempo de vuelo}
\newacronym{PDP}{PDP}{perfil de potencia en función del retardo}
\newacronym{RMSDS}{RMS-DS}{dispersión cuadrática media del retardo}
\newacronym{SNR}{SNR}{relación señal a ruido}
\newacronym{SSNR}{SSNR}{relación señal a ruido para sensado}
\newacronym{CFO}{CFO}{desajuste de frecuencia de portadora}
\newacronym{SFO}{SFO}{desajuste de frecuencia de muestreo}
\newacronym{STO}{STO}{desajuste temporal de muestreo}
\newacronym{CPO}{CPO}{desajuste de fase de portadora}
\newacronym{PEOC}{PEOC}{calibración que ignora el error de fase}
\newacronym{UPEC}{UPEC}{calibración con error de fase uniforme}
\newacronym{PEAC}{PEAC}{calibración que considera el error de fase}
\newacronym{PHY}{PHY}{capa física}
\newacronym{CSS}{CSS}{espectro ensanchado mediante chirridos}
\newacronym{FFT}{FFT}{transformada rápida de Fourier}
\newacronym{IFFT}{IFFT}{transformada rápida inversa de Fourier}
\newacronym{FIM}{FIM}{matriz de información de Fisher}
\newacronym{CRLB}{CRLB}{cota inferior de Cramér--Rao}
\newacronym{LOS}{LOS}{línea de vista}
\newacronym{NLOS}{NLOS}{sin línea de vista}
\newacronym{AWGN}{AWGN}{ruido gaussiano blanco aditivo}
\newacronym{LTI}{LTI}{lineal e invariante en el tiempo}
\newacronym{LTV}{LTV}{lineal y variante en el tiempo}
\newacronym{WSSUS}{WSSUS}{estacionario en sentido amplio con dispersión no correlacionada}
\newacronym{CP}{CP}{prefijo cíclico}
\newacronym{PSD}{PSD}{densidad espectral de potencia}

% ML / Representation Learning
\newacronym{DL}{DL}{aprendizaje profundo}
\newacronym{ML}{ML}{aprendizaje automático}
\newacronym{SSL}{SSL}{aprendizaje autosupervisado}
\newacronym{JEPA}{JEPA}{Joint-Embedding Predictive Architecture}
\newacronym{PINN}{PINN}{red neuronal informada por la física}
\newacronym{BSS}{BSS}{separación ciega de fuentes}
\newacronym{PCA}{PCA}{análisis de componentes principales}
\newacronym{ICA}{ICA}{análisis de componentes independientes}
\newacronym{NMF}{NMF}{factorización matricial no negativa}
\newacronym{NMSE}{NMSE}{error cuadrático medio normalizado}
\newacronym{MAE}{MAE}{error absoluto medio}
\newacronym{RMSE}{RMSE}{raíz del error cuadrático medio}
\newacronym{OOD}{OOD}{fuera de distribución}
\newacronym{SMPL}{SMPL}{Skinned Multi-Person Linear Model}

% Health / Human sensing
\newacronym{RR}{RR}{frecuencia respiratoria}
\newacronym{HR}{HR}{frecuencia cardiaca}
\newacronym{ECG}{ECG}{electrocardiograma}
\newacronym{PPG}{PPG}{fotopletismograma}
\newacronym{GT}{GT}{valor de referencia}

% ======================================================
%                COMANDOS ÚTILES
% ======================================================
\newcommand{\cpx}{\mathbb{C}}
\newcommand{\real}{\mathbb{R}}
\newcommand{\E}{\mathbb{E}}
\newcommand{\e}{\mathrm{e}}
\newcommand{\jj}{\mathrm{j}}
\newcommand{\calP}{\mathcal{P}}
\newcommand{\calO}{\mathcal{O}}
\newcommand{\calI}{\mathcal{I}}
\newcommand{\calL}{\mathcal{L}}
\newcommand{\calD}{\mathcal{D}}
\newcommand{\calR}{\mathcal{R}}
\newcommand{\vect}[1]{\bm{#1}}
\newcommand{\mat}[1]{\bm{#1}}
\DeclareMathOperator{\wrap}{wrap}
\DeclareMathOperator{\diag}{diag}
\DeclareMathOperator{\rank}{rank}
\DeclareMathOperator{\tr}{tr}

% ======================================================
%                    DOCUMENTO
% ======================================================

% Evita estirar verticalmente las páginas con tablas extensas.
\raggedbottom
